\begin{textsample}{POS Dim 1 – persona\_grok – Score 21.00 – t450\_grok.txt}  \label{ex:f1_pos_032}
In 2019, massive climate \textbf{marches} brought millions of \textbf{people} together around the \textbf{world}, amplifying urgent calls for action on the climate crisis.
These gatherings showcased the power of collective voices filling the \textbf{streets}, from \textbf{school} strikers to global protests demanding systemic change.
But the global pandemic has changed \textbf{everything}, enforcing physical distancing measures that have limited traditional \textbf{street} protests and large-scale \textbf{demonstrations}.
The need to protect public health has forced a pause on the \textbf{kind} of in-person mobilizations that defined the \textbf{movement} just a \textbf{year} ago.
Yet, this hasn't stopped \textbf{activists} worldwide from demanding climate action.
Instead, they 've adapted with \textbf{creativity} and determination, finding innovative \textbf{ways} to keep the momentum alive.
Online actions have surged, connecting \textbf{people} across borders through digital \textbf{platforms}.
Hashtags like \#DigitalStrike, led by Fridays For Future, have trended globally, rallying support and raising \textbf{awareness} from home.
Virtual \textbf{marches} have brought the energy of protests into the digital realm, while holographic projections have created striking visual displays without crowds.
In some places, distanced physical protests maintain the \textbf{spirit} of assembly while adhering to safety guidelines.
Social \textbf{media} has become a canvas for solidarity symbols, with users \textbf{sharing} \textbf{images} and messages that unite the \textbf{movement}.
Even from isolation, \textbf{activists} are projecting banners from their homes using lights and shadows, turning personal spaces into public statements.
These examples inspire \textbf{us} all, showing how the climate \textbf{movement} is evolving in the face of adversity.
Isolation won't silence our voices—it 's only making them more resilient and inventive.
As we navigate this crisis, \textbf{let} 's turn it into an opportunity to build a \textbf{better} post-COVID \textbf{world}, \textbf{one} that prioritizes the health of \textbf{people} and the planet through bold, equitable climate solutions.

% matched lemmas: activist, awareness, creativity, demonstration, everything, good, image, kind, let, march, medium, movement, one, people, platform, school, share, spirit, street, us, way, world, year
\end{textsample}
